\section{Graphs, homology, and the surface table}
\label{sec:algtop}
The motivating article singles out the seven bridges, Euler's formula, and a
table of invariants of closed surfaces. These are three different kinds of
claim. A walk is a finite word, a homology computation is an identity about
integer or binary matrices, and the identification of an arbitrary space
with a particular surface requires a representation theorem. We prove the
first two directly. For surfaces we specify the presentations on which the
computations are performed.

\subsection{Paths and bridges}
A finite multigraph has vertices $0,\ldots,v-1$ and numbered edges
$0,\ldots,e-1$ with two specified ends each. The two ends may coincide.
An oriented edge is an edge number and a bit choosing its direction;
reversal flips that bit, even on a loop. A path stores its initial vertex
and a composable list of oriented edges. Degree counts ends with
multiplicity. Starting at a vertex, mark it, then repeatedly mark a new
vertex attached to a marked one, storing the attaching edge as its parent.
Stop at exhaustion or after $v$ passes, with one new mark per successful
pass. An edge cannot then lead from a marked to an unmarked vertex.
Induction along any given path proves it stays in the marked set. Parent
pointers, read in reverse marking order, give paths of length $<v$ to all
marks. Repeating from the least unmarked vertex gives a spanning forest
with exactly $v-c$ edges, where $c$ is the number of roots (components).
It has no cycle, since the last-marked vertex on a cycle would need two
earlier neighbors along forest edges. This argument counts a doubled edge
or a loop as a cycle when applicable.

\begin{theorem}[Euler trails and the seven bridges]\label{idx-algebraic29}
For a graph with at least one edge, an edge-once trail exists precisely
when all positive-degree vertices are in one component and the number of
odd-degree vertices is zero or two. In the former case the trail closes.
In the latter its endpoints are exactly the two odd vertices.
\end{theorem}
\begin{proof}
A trail connects every visited vertex to its start. Pair each arrival with
the next departure at an internal visit. All ends are paired except the
initial departure and terminal arrival, which also pair if the endpoints
coincide. This proves necessity. Suppose first all degrees are even.
Starting at an incident vertex, traverse the least unused incident edge
until there is none. Stopping anywhere other than the start would leave an
odd number of used ends at a vertex with no unused ends. Thus we have a
closed trail. Removing it preserves even degrees. If some edge is unused,
some visited vertex has an unused incident edge: otherwise no edge could
leave the visited set, and connectedness of the support would force every
incident vertex to have been visited and every edge used. Make another
closed trail from such a vertex in the residual graph, and splice it into
the first. Each splice uses a fresh edge. At most $e$ walks of length at
most $e$ suffice. With two odd vertices first join them by an extra edge.
Cut the resulting cyclic list at this edge and drop it. For the historical
bridge graph the four degrees are $5,3,3,3$ (one vertex has two edges to
each of two banks and one to the other island; each bank has one further
edge to that island). Its seven bridges therefore admit no such route.
\end{proof}

\begin{theorem}[The edge-path group of a graph]\label{idx-algebraic56}
At a root in a connected nonempty graph the group of based edge paths
modulo insertion and deletion of $a\bar a$ is the free group on
$e-v+1$ generators. Equality has a primitive-recursive normal form.
\end{theorem}
\begin{proof}
Read a word on signed letters with a stack, popping if the incoming letter
inverts the top and pushing otherwise. Its stack is reduced after every
step. Reading $a\bar a$ restores any reduced stack: if $a$ pops its
inverse, the exposed letter cannot be $a$, because the old stack was
reduced. In the other case the push is immediately undone. Thus insertion
and deletion preserve the answer in every context. Conversely the scanned
prefix is equivalent to its current stack, by induction on the scan. The scan records at most one adjacent-pair deletion per input letter,
so it also produces an explicit move transcript. This
proves the normal form theorem for free words (and the same argument
respects endpoints for edge paths).

Fix the parent tree and let $t_x$ be its route from the root to $x$.
Send a based path to its list of non-tree edges, oriented relative to
their given directions, then reduce. A generator corresponding to
$a:x\to y$ lifts to $t_xa\bar t_y$. Its inverse lifts to the reverse
loop. Projecting this lift gives the generator. Conversely insert
$\bar t_x t_x$ between consecutive edges at $x$ (each detour is obtained by nested
inverse-pair insertions) to write a based path as
a product of the lifted single-edge loops. For a tree edge its loop
cancels completely, since one endpoint is the parent of the other.
The constructions are inverse and concatenate modulo cancellation, which
also proves the group laws from the word operations.
Lifting a word of length $l$ uses at most $(2v+1)l$ edge symbols.
This is the combinatorial edge-path interpretation of the graph's
fundamental group; its definition here quantifies over finite paths and
moves.
\end{proof}

\subsection{Finite chain complexes and Betti numbers}
A simplicial complex on $v$ available labels is a table on the $2^v$
subsets, with the empty subset present and every subset of a present face
present. Labels absent as singletons are harmless. Order faces by their
bit masks. In degree $k\ge0$ the binary chain group $C_k$ consists of
masks on the $f_k$ faces of size $k+1$, with exclusive-or as addition.
In ordinary chains $d_0=0$. In positive degree, $d_k$ sends each face to
the sum of its one-vertex deletions and extends by addition. Write
$Z_k=\ker d_k$, $B_k=\operatorname{im}d_{k+1}$. For $k\ge2$ each
two-vertex deletion occurs twice in $d_{k-1}d_k$; for $k=1$ the
composite vanishes by $d_0=0$. Thus $B_k\subseteq Z_k$. Two cycles
represent the same homology class exactly when their sum is a boundary.

\begin{lemma}[Binary elimination with witnesses]\label{idx-algebraic103}
For an $m$ by $n$ binary matrix $A$, a bounded algorithm produces $r$
independent columns spanning the image, and $n-r$ independent kernel
vectors spanning the kernel. It also solves $Ax=y$ or detects failure.
\end{lemma}
\begin{proof}
Process columns in order. Store each surviving reduced column together
with the domain vector mapping to it. Give it the highest occupied row
as pivot. For the next column start with its domain unit vector; scan rows
from high to low, eliminating occupied old pivots by adding the stored
pair. If a nonzero residual has a new highest pivot, store it with its
domain vector. If the residual is zero, store its domain vector as a
kernel vector. (Eliminating all old pivots in one downward pass preserves
every cancellation above the row being processed.) The domain vectors
recorded after all columns, whether residuals are zero or nonzero, form a
basis: the vector for column $j$ contains bit $j$ and uses no later bit.
Descending elimination hence expresses any vector uniquely in this basis.
The nonzero residuals are independent because a largest pivot in a
nonempty sum survives. Consequently the zero-residual domain vectors span
exactly the kernel and the nonzero residuals span the image. To solve,
reduce $y$ in descending pivot order while accumulating the same domain
vectors. A nonzero remainder cannot be in the span because it contains no
pivot; a zero remainder exhibits a solution. If original independent
columns are desired, take the original column indices with nonzero
residual: at each such step subtracting a combination of earlier
residuals preserves the generated span and enlarges its dimension by one.
\end{proof}

\begin{theorem}[Euler--Poincar\'e and cohomology]\label{idx-algebraic131}
With $r_k=\operatorname{rank}d_k$, the ordinary homology classes have
unique binary coordinates of length
$\beta_k=f_k-r_k-r_{k+1}$. The same dimension is obtained for the
cohomology of the transposed boundary matrices. In particular,
\[
 \chi(K):=\sum_{k\ge0}(-1)^k f_k
        =\sum_{k\ge0}(-1)^k\beta_k .
\]
For subcomplexes $K,L$ on the same labels,
$\chi(K\cup L)+\chi(K\cap L)=\chi(K)+\chi(L)$.
\end{theorem}
\begin{proof}
Start a basis with the independent boundary vectors in $C_k$. Append
successively the computed kernel generators whose reduction is nonzero.
Keep the residuals and pivots as in the lemma. The final list spans the
kernel. Each cycle has unique coordinates in it, by descending pivot
elimination; two cycles differ by a boundary just when the appended
coordinates agree. A binary space with bases of sizes $a,b$ has both
$2^a$ and $2^b$ elements by unique coordinates, hence $a=b$ (alternatively
use elimination). The final list thus has length $f_k-r_k$. This yields
the claimed number of appended coordinates. For the transpose, the
domain change of basis in the elimination lemma puts $A$ in the form
consisting of independent image vectors and zeros. Extend these vectors
to a codomain basis by appending independent unit vectors. In these two
bases $A$ has exactly $r$ diagonal ones and remaining entries zero, up to
a permutation. Transposing the elementary invertible basis changes
preserves their invertibility, proving equality of ranks. Apply the same
kernel-quotient count to the transposed complex. In the alternating sum
the two appearances of each nonzero boundary rank cancel, including the
end cases since $d_0=0$ and degrees above $v$ vanish.
Inclusion--exclusion follows by testing membership of each nonempty face
in the two tables.
\end{proof}

\begin{corollary}[Graphs and degree zero]\label{idx-algebraic166}
For a finite complex $\beta_0$ is the number $c$ of components of its
vertex-edge graph. For a finite graph, allowing loops and repeated edges
by using their binary incidence columns, $\beta_1=e-v+c$ and
$\chi=v-e$.
\end{corollary}
\begin{proof}
Each edge boundary has even total parity in each component. Conversely
the boundary of a parent route from $x$ to its root is the pair of
endpoints (zero when they coincide). Summing these routes over the marked
bits of a vertex chain with even parity in each component reproduces that
chain as a boundary. The component parities are independent since each
root is a representative. Thus $r_1=v-c$. A graph has $d_2=0$, proving
the degree-one formula.
\end{proof}

\subsection{Cones, spheres, and invariance under subdivision}
Use the augmented complex temporarily, adjoining one empty face in degree
$-1$, with $d$ of a vertex the empty face and $d$ of the empty face zero.
This is an index shift of a finite list of matrices. For any simplicial
cone with apex $u$ (that is, $s\cup\{u\}$ is present for every face
$s$), let $h(s)$ adjoin $u$ if absent and be zero if present. If $u$ is
absent, its deletion from $h(s)$ gives $s$, and deleting any other vertex
matches a summand in $hd(s)$. If $u$ is present, only its deletion in
$hd(s)$ contributes. This includes the empty face. Hence
$dh+hd=1$. In particular a reduced cycle $z$ has filling $h(z)$.

\begin{theorem}[Simplex and boundary sphere]\label{idx-algebraic193}
The full simplex on $n+1$ vertices, $n\ge0$, has vanishing reduced
binary homology. For $n\ge1$, its boundary has reduced homology of
dimension one in degree $n-1$ and zero elsewhere. The top cycle is the
sum of its facets. In particular a triangle boundary has ordinary Betti
numbers $(1,1)$, and a tetrahedral boundary has $(1,0,1)$ and
$V-E+F=2$.
\end{theorem}
\begin{proof}
The full simplex is a cone. A reduced boundary-sphere cycle $z$ of
degree $k<n-1$ fills by the cone formula in the full simplex; its filling
still has degree below $n$, so uses only proper faces. In degree $n-1$
the cone filling is a multiple of the unique $n$-simplex, so every
cycle is a multiple of its boundary. That boundary has distinct facets,
is nonzero, and has no possible filling among the proper faces. In
higher degrees there are no chains. For a nonempty complex, sorting a
vertex chain by its parity splits off the class of a chosen vertex in
ordinary degree zero, leaving reduced degree zero; positive degrees are
unaffected.
\end{proof}

A simplicial vertex map sends a face to a face as an unordered image. On
chains let $F$ send a basis face to its image if the vertices remain
distinct and to zero otherwise. It sends the empty face to itself.
It commutes with $d$: for an injectively mapped face, deletions match
bijectively. If the image has deficit one, only the two deletions from
the duplicated pair give nonzero terms, and their images agree. With
larger deficit every deletion still has a repetition. This rule
respects identity and composition since repeated images cannot regain
distinctness under a second map. A chain map therefore sends cycles to
cycles and boundaries to boundaries, inducing a map on homology. The
algorithm above supplies its matrix by reducing the images of class
representatives.

\begin{lemma}[A cone filling for chain homotopies]\label{idx-algebraic227}
Suppose augmentation-preserving chain maps $F,G$ from a simplicial
complex to another are carried by cones $Q(s)$: both images of each
nonempty basis face $s$ are supported there and
$Q(t)\subseteq Q(s)$ for nonempty $t\subseteq s$. Then there is a
computed family of binary matrices $P$ of degree $+1$ such that
$dP+Pd=F+G$. In particular $F,G$ induce the same map on homology.
\end{lemma}
\begin{proof}
Put $P(\varnothing)=0$. In increasing face size set
\[
 P(s)=h_{Q(s)}\bigl(F(s)+G(s)+P(d s)\bigr).
\]
The chain in parentheses lies in $Q(s)$ by the containment condition and
the previous steps. Its boundary is zero: use the two chain map
identities and $dP(d s)=(F+G)(d s)+P(d^2s)$, already established on
proper faces (including the empty face). Hence its cone filling has the
asserted boundary. Extending additively proves the identity. On an
ordinary positive-degree cycle the difference of images is the boundary
of $P(z)$; in degree zero the same conclusion holds since
$P(\varnothing)=0$.
\end{proof}

\begin{theorem}[Contiguity, elementary collapse, and subdivision]\label{idx-algebraic250}
Contiguous simplicial maps induce equal maps on homology. Removing a
maximal simplex $s$ together with a nonempty codimension-one face $t$
having no other proper coface preserves homology. Barycentric subdivision
induces a computable isomorphism on homology in all degrees.
\end{theorem}
\begin{proof}
Contiguity means that the union of the two images of each face is a face.
Take $Q(s)$ to be the full simplex on that union and apply the cone
lemma. Consequently a finite chain of contiguities gives the same map;
mutual inverse maps up to such chains induce inverse homology maps.

For a collapse write $d s=t+b$, where $b$ uses only remaining faces.
The projection to the remaining chains fixes these, sends $s$ to zero,
and $t$ to $b$. It commutes with $d$ since $d b=d t$ and $t$ occurs
in the boundary of no other simplex. Its composite with inclusion is
the identity on remaining chains. On the original chains the difference
from the identity is $dP+Pd$ with $P(t)=s$ and $P=0$ elsewhere, as
one sees on $s,t$ and each remaining face. This also proves invariance
for a supplied finite sequence of collapses or their inverses.

Vertices of $\operatorname{sd}K$ are the nonempty faces of $K$, and
its simplices are strictly increasing chains of such faces. On augmented
chains define $S(\varnothing)=\varnothing$ (the basis vector).
Given a nonempty face $s$, cone $S(d s)$ to the vertex $s$ in
$\operatorname{sd}(2^s)$, the subdivision of the full simplex on $s$.
This complex is a cone since every chain can be extended by $s$ unless
it ends there. By induction $S(d s)$ is a cycle in that cone and its
filling has boundary $S(d s)$, so $S$ is a chain map with the indicated
support. Map a subdivision vertex $t$ to the greatest original vertex
of $t$. This assignment is simplicial since all images of a chain
belong to its largest face. Call its chain map $L$. The maps $LS$ and
the identity on $K$ are carried on $s$ by the full simplex on $s$.
The maps $SL$ and the identity on $\operatorname{sd}K$ are carried on
a nonempty chain ending in $s$ by $\operatorname{sd}(2^s)$.
These carriers are nested under taking faces. Apply the lemma twice to
obtain inverse homology maps. Iteration along any supplied finite list of
subdivisions or simplicial isomorphisms preserves the calculated groups
and Euler characteristic.
\end{proof}

\subsection{Euler's formula for a rational convex polyhedron}
A vertex presentation lists distinct points $p_1,\dots,p_m$ in
$\mathbb Q^3$ of affine dimension three. With each $p$ it gives a
rational covector $u_p$ with $u_p(q-p)>0$ for all other points $q$.
This certificate ensures that every listed point is a vertex. The
polyhedron means their convex hull. Its facets are the maximal lists
in a supporting plane spanning that plane, and its edges are its
exposed segments. These data can all be tested with rational arithmetic.
In particular supporting facet candidates are enumerated by triples.
Here convex combinations can be read throughout as rational ones.

\begin{lemma}[Planar sections of a finite hull]\label{idx-algebraic302}
For noncollinear rational points in a plane, their hull has a boundary
polygon. A rational affine form $x$ taking distinct values at the points
splits its boundary into two chains from the least to the greatest
$x$, along which $x$ increases. Consequently a cut $x<a$ avoiding the
polygon vertices selects one consecutive cyclic block of vertices, unless
it selects all vertices or none.
\end{lemma}
\begin{proof}
Use $x$ as the first affine coordinate and take an independent rational
coordinate $y$. At each rational abscissa $t$ between the extremes, the
possible ordinates form the interval between the lowest and highest
ordinates obtained by interpolating pairs of the given points spanning
$t$ (include a point at $t$). Indeed, in a convex combination with
three positive coefficients there is a nonzero rational change of those
three coefficients preserving their sum and their weighted sum of $x$.
It has both signs. Choose its sign so that the weighted $y$ does not
decrease, and move until a coefficient first becomes zero. Iterate at
most the size of the support to get an upper bound from a pair. For a
lower bound reverse the choice of sign. The extreme ordinates are
attained and all intermediate ordinates are mixtures of the two.
The resulting upper envelope is concave by taking mixtures of its
maximizers. Between successive input abscissae the candidate segments
are fixed, so the upper envelope is also convex there as a maximum of
affine functions. It is therefore affine on each such interval.
There is no jump at an interior input abscissa, by the concavity
inequality using points on both sides; at either extreme the finitely
many candidate segments reaching it share the unique endpoint.
The lower envelope has the reversed argument. Noncollinearity makes
the two envelopes distinct at every interior abscissa: an interior
triangle has positive vertical thickness at some abscissa, and mixtures
with either endpoint propagate positive thickness throughout the
interior. Removing collinear intermediate points leaves two strictly
$x$-monotone polygonal chains, meeting only at their ends. Each corner
is an input point: an extremizer at a corner supported by two different
abscissae would, by concavity (or convexity), force the same affine
segment on both sides. Each resulting side has a supporting line, as
follows by monotonicity of the slopes of a concave or convex function.
These two chains give the asserted cyclic block rule.
\end{proof}

\begin{theorem}[Euler's polyhedron formula]\label{idx-algebraic343}
For a full-dimensional rational convex vertex presentation,
$V-E+F=2$, with polygonal faces counted before any triangulation.
\end{theorem}
\begin{proof}
At $p$ form the planar hull $Q_p$ of the normalized directions
$q'=(q-p)/u_p(q-p)$, lying in $u_p=1$. It is two-dimensional since the
rays span three-space. Its points $q'$ are distinct: two on the same
ray would make the nearer input vertex a strict convex combination of
$p$ and the farther one, contradicting its exposing certificate.
Every affine form on $u_p=1$ extends to a linear form on the ambient
vector space (extend its constant using $u_p$). Thus a supporting line
of $Q_p$ gives a supporting plane of the original hull through $p$.
A corner of $Q_p$ is exposed by the sum of inward slacks of its two
incident sides, giving exactly one edge from $p$. Conversely an exposed
edge from $p$ must expose just one of the normalized directions, a corner.
A side of $Q_p$ gives exactly one facet through $p$: its two extreme
directions span a plane with $p$, and all its directions give exactly
the input vertices on that plane. Each facet through $p$ gives a side
in reverse. This proves the incidence correspondence used below.

Choose a rational linear height $h$ separating the input vertices and
also separating all pairs of normalized directions within each $Q_p$.
Testing $(1,j,j^2)$ for $j=0,1,\ldots,2M$ suffices for $M$ pairwise
inequalities to avoid: each nonzero difference gives a nonzero polynomial
of degree at most two with at most two integer roots. Heights of the
normalized directions are nonzero, with the sign of $h(q)-h(p)$.
At $p$ count $1$ for the vertex, subtract $1$ for each edge pointing
downward from $p$, and add $1$ for each facet whose greatest vertex is
$p$. The latter facets correspond to sides of $Q_p$ with both ends
negative, because every direction on the side is a mixture of its ends.
These contributions count each edge and each facet exactly once globally.

If all corners are positive, all normalized directions are positive by
convexity and $p$ is the global minimum. Its contribution is $1$.
If all are negative, $p$ is the global maximum and its contribution is
$1-k+k=1$, for $k$ corners. In the mixed case let $k$ be the number of
negative corners. They form one consecutive proper block by the planar
lemma for $x=h,a=0$, and exactly $k-1$ sides have both ends in it. The
contribution $1-k+(k-1)$ vanishes. There is one global minimum and one
global maximum, proving the formula. The planar argument uses at most
all pairs of points per envelope ordinate. Affineness reduces each
inequality about an interpolated interval to its endpoints. Thus the
sweep needs only bounded lists and rational operations in iPRA.
\end{proof}

\subsection{The integer homology and orientability of the polygon models}
Work now over $\mathbb Z$. Use $g$ handles and $c$ cross-caps in the
standard polygon word
\[
 w=\prod_{i=1}^{g}a_ib_ia_i^{-1}b_i^{-1}
       \prod_{j=1}^{c}y_jy_j .
\]
The formula below concerns the polygon with paired edges indicated by
this word, including the orientable models $c=0$ and the nonorientable
ones $c>0$. The empty word is handled by the tetrahedron boundary.
No recognition or classification algorithm for an arbitrary supplied
surface is needed.

\begin{theorem}[The surface table over the integers]\label{idx-algebraic402}
Write $m=2g+c$. The connected polygon model has Euler characteristic
$2-m$. Its homology in degrees $0,1,2$ is respectively
\[
\begin{array}{c|ccc}
 &H_0&H_1&H_2\\ \hline
 c=0 &\mathbb Z&\mathbb Z^{2g}&\mathbb Z\\
 c>0 &\mathbb Z&\mathbb Z^{m-1}\oplus\mathbb Z/2&0 .
\end{array}
\]
It is orientable exactly when $c=0$. In particular the sphere, torus,
and double torus have $(b_0,b_1,b_2)=(1,0,1),(1,2,1),(1,4,1)$.
The projective plane and a Klein bottle model have $(1,0,0)$ and
$(1,1,0)$ respectively, each with one order-two summand in degree one.
\end{theorem}
\begin{proof}
Suppose the word has length $l>0$. All boundary vertices are identified:
in a commutator block with successive corners $0,1,2,3,4$ the first
pair gives $0=3,1=2$ and the second gives $1=4,2=3$. In a square
block the two copies identify all three corners. Blocks share endpoints.
Call the common vertex $p$, add a center $o$, and draw $l$ spokes
$r_i:o\to p$ to successive corners ($i$ modulo $l$). These are
distinct numbered edges. Let $e_j$ denote the oriented edge for each of
the $m$ distinct boundary letters. The oriented triangle between spokes
$i,i+1$ has boundary
\[
 dT_i=r_i+s_ie_{j(i)}-r_{i+1},\qquad
 dr_i=p-o,\quad de_j=0 ,
\]
where $s_i$ is the sign of the letter. This is the ordinary oriented
triangle-edge-vertex incidence complex of the glued fan, with face
identifications allowed. There are $2$ vertices, $m+l$ edges and $l$
triangles. This proves the Euler count and $H_0=\mathbb Z$.

A $1$-cycle has total spoke coefficient zero, so its unique coordinates
are the $e_j$ and $q_i=r_i-r_0$ for $1\le i<l$. Put $q_0=q_l=0$.
The boundary relations are $q_{i+1}=q_i+s_ie_{j(i)}$. Successive
substitution eliminates all $q_i$ leaving just
$u=\sum_i s_i e_{j(i)}=0$. Explicitly, replace each $q_i$ by the prefix
sum up to index $i-1$, using the first $i$ triangle boundaries.
These substitutions kill the first $l-1$ relations and send the last
one to $u$; conversely $u=\sum_i dT_i$. Thus a cycle is a boundary
exactly when its final $e$-vector is an integer multiple of $u$.
The spoke coefficients of $d\sum_i t_iT_i$ vanish exactly when
$t_i=t_{i-1}$ for every $i$. The remaining boundary is then $t_0u$.
There are no $3$-chains. For handles alone $u=0$, giving the first
row. With cross-caps, $u=2(y_1+\cdots+y_c)$, where letters now denote
basis vectors. Replacing $y_1$ by $y_1+\cdots+y_c$ is an invertible
integer basis change with inverse obtained by subtraction. The new
coordinate is reduced modulo two, proving its exact order and both
homology formulas. A consistent orientation chooses $t_i\in\{1,-1\}$
for every triangle with paired edge incidences cancelling. The spokes
force a single sign throughout, which succeeds precisely when the
paired boundary incidences have opposite directions, i.e. $c=0$.

For the empty word use the tetrahedral boundary with oriented faces
and alternating deletion signs. In the full simplex, prefixing the
least vertex if absent (and returning zero otherwise) gives
$dh+hd=1$ on augmented integer chains: deletion of the prefixed vertex
has positive sign, and the signs of all later deletions reverse. When
the least vertex is already present only its deletion contributes to
$hd$. The preceding simplex-boundary proof now works over $\mathbb Z$,
giving the sphere row. Its signed facet boundary supplies a consistent
orientation. Finally the usual Klein bottle polygon can also be read as
$a b a^{-1} b$; its paired corners are all identified and its exponent
vector is $2b$, so the same spoke calculation proves its stated row
directly.
\end{proof}

\subsection{Finite filtrations and their barcodes}
Let $K_0\subseteq\cdots\subseteq K_N$ be finite face tables on common
labels. In any fixed degree, the preceding algorithms turn the inclusions
into binary matrices $A_i:V_i\to V_{i+1}$ on ordinary homology
coordinates. Examples include the proximity complexes of a finite
rational metric space, whose face test at threshold $t$ is
$d(x,y)\le t$ for every pair in the face. The following applies to
every finite sequence of binary matrices.

\begin{theorem}[Interval decomposition and recovery from ranks]\label{idx-algebraic480}
There are explicit bases in which the sequence is a direct sum of
intervals $[b,d]$, $0\le b\le d\le N$: each interval contributes one
coordinate at each covered index and identity maps between these
coordinates. If $R(i,j)$ is the rank of the composite $V_i\to V_j$
for $i\le j$ (identity for $i=j$), the multiplicity of $[b,d]$ is
\[
 R(b,d)-R(b-1,d)-R(b,d+1)+R(b-1,d+1)
\]
with out-of-range ranks read as zero. Thus the multiset of intervals
and every persistent Betti number are computable.
\end{theorem}
\begin{proof}
If there is any nonzero space, take the earliest one, $V_b$, and a unit
vector $x_b$. Propagate it forward until it becomes zero or the sequence
ends, and let $x_d$ be its last nonzero image. These lines form an
interval submodule since the outgoing map kills $x_d$ if it exists.
Choose a coordinate functional $\lambda_d$ with
$\lambda_d(x_d)=1$. Pull it back along each composite from $t$ to $d$,
giving $\lambda_t(x_t)=1$ for $b\le t\le d$. The projections
$z\mapsto\lambda_t(z)x_t$ there and zero elsewhere commute with the
arrows: in the interior this is the pullback identity, after $d$ the
line dies, and the space before $b$ is zero. Thus their kernels form
a complementary sequence. Elimination finds bases of these kernels;
every vector splits uniquely by subtracting its projection. Apply
the operation to the complement. Total dimension drops by at least one
each time, so at most $\sum_i\dim V_i$ steps give an interval basis.
In this basis $R(i,j)$ counts intervals starting at or before $i$ and
ending at or after $j$. Taking successive differences retains precisely
the intervals starting at $b$ and ending at $d$. This proves uniqueness
of the multiplicities and the claimed formula.
\end{proof}

\subsection{Gluing and reconstruction from a cover}
For the following exactness statement reduced groups use the augmented
chain convention. Even the complex having only the empty face has its
degree $-1$ chain group. Ordinary groups work as well when the groups
in negative degrees are set to zero.

\begin{theorem}[Mayer--Vietoris with a computed connecting map]\label{idx-algebraic519}
Let $K=A\cup B$, with $A,B$ finite subcomplexes, and $I=A\cap B$.
For binary coefficients the following sequence is exact in every degree:
\[
 H_k(I)\xrightarrow{u} H_k(A)\oplus H_k(B)
 \xrightarrow{v}H_k(K)\xrightarrow{D}H_{k-1}(I)
 \xrightarrow{u} H_{k-1}(A)\oplus H_{k-1}(B).
\]
Here $u[z]=([z],[z])$, $v([a],[b])=[a+b]$.
\end{theorem}
\begin{proof}
Split a cycle $z=a+b$ by assigning each face in it to $A$ if possible
and to $B$ otherwise. Then $da=db$ is supported in $I$ and is a cycle.
Put $D[z]=[da]$. A second split differs by a pair $(i,i)$ in $I$,
so its derivative changes by $di$. A change $z\mapsto z+dw$ permits
the split $(a+dp,b+dq)$ using any split $w=p+q$, which leaves $da$
unchanged. Thus $D$ is well-defined and additive, the latter by taking
the sum of any splits.

At $H_k(A)\oplus H_k(B)$ a diagonal class sums to zero. Conversely if
$a+b=dw$, split $w=p+q$ and obtain $a+dp=b+dq$, a common cycle in
$I$ giving the desired preimage. At $H_k(K)$ a pair of cycles supplies
a split with zero derivative. Conversely if $da=di$ in $I$, then
$a+i,b+i$ are cycles summing to $z$. At $H_{k-1}(I)$ the class of
$da=db$ vanishes in both sides. Conversely if a cycle $i$ in $I$
bounds $a$ in $A$ and $b$ in $B$, then $a+b$ is a cycle in $K$ with
$D[a+b]=[i]$. The elimination lemma recovers each filling in these
tests and hence a preimage coordinate vector at every exact spot.
\end{proof}

\begin{theorem}[Finite nerve lemma for cone intersections]\label{idx-algebraic549}
Let $K$ be the union of a finite list of subcomplexes $K_i$. Discard
members without vertices. For every nonempty set $J$ of indices whose
intersection $K_J=\bigcap_{j\in J}K_j$ has a vertex, suppose $K_J$
is a simplicial cone. The nerve $N$ has such sets $J$ as nonempty
faces. There are inverse computable isomorphisms between the binary
homology groups of $K$ and $N$.
\end{theorem}
\begin{proof}
Choose the least apex $a_J$ passing the cone test on $K_J$.
There is a simplicial map $F:\operatorname{sd}K\to\operatorname{sd}N$
sending $s$ to $I(s)=\{i:s\in K_i\}$. It reverses inclusions
and sends a chain of faces to a possibly repeated chain. The set $I(s)$
is a face of $N$ since $s$ has a vertex. In the other direction
$G:\operatorname{sd}N\to K$ sends $J$ to $a_J$. To see that it is
simplicial, start with the apex at the end of a chain
$J_0\subsetneq\cdots\subsetneq J_k$, and adjoin the earlier apices
in descending order. At step $i$ the later face belongs to
$K_{J_{i+1}}\subseteq K_{J_i}$, where it can be enlarged by $a_{J_i}$.
Thus the entire image belongs to $K_{J_0}$.

Use $S$ and $L$ for the subdivision and last-vertex chain maps of $K$.
On a face of $\operatorname{sd}K$ ending in $s$, both $GF$ and
$L$ have support in $K_{I(s)}$. The former assertion uses that
$I(s)$ is the smallest index set in the image chain, while the latter
uses $s\in K_i$ for all $i\in I(s)$. These cones nest under passage to
a subface, so $GF$ and $L$ act identically on homology. It follows
that $GFS$ acts as the identity. For a face of $\operatorname{sd}N$
beginning in $J_0$, every original face $s$ occurring as a subdivision
vertex in $SG$ belongs to $K_{J_0}$. Thus every vertex $I(s)$ of its
$F$-image contains $J_0$. The subcomplex of $\operatorname{sd}N$
consisting of chains of supersets of $J_0$ is a cone with apex $J_0$,
and contains the original face too. When the least index set grows
this cone shrinks. The cone homotopy lemma gives $FSG=1$ on homology.
Consequently $G$ and $FS$ induce inverse isomorphisms between the
homology of $K$ and of $\operatorname{sd}N$. Apply the subdivision
theorem to $N$. The case with no vertices is covered by the unique
maps on the complexes having only the empty face.
\end{proof}

\begin{corollary}[Dual complexes of a finite relation]\label{idx-algebraic589}
For a binary relation $R\subseteq X\times Y$ on finite lists, give a
nonempty subset of $X$ the face test that its elements have a common
$R$-neighbor in $Y$, and use the transposed test on $Y$.
The two complexes have isomorphic binary homology in every degree.
\end{corollary}
\begin{proof}
Cover the first complex by the simplices on the neighborhoods of the
individual elements of $Y$. Every intersection with vertices is a full
simplex. Its nerve is exactly the second complex, with the isolated
unused labels omitted. Apply the finite nerve lemma.
\end{proof}

\subsection{Arithmetic and Lean implementation details}
All spaces of chains here have explicit carriers. On $v$ possible
vertices a face is below $2^v$ and a table of faces below $2^{2^v}$.
The group $C_k$ uses bit masks below $2^{f_k}$. An $m$ by $n$ binary
matrix fits in $mn$ bits. A face of a barycentric subdivision is itself
a subset of the fewer than $2^v$ nonempty faces. Enumerate subsets and
test nestedness by comparing all pairs. A next subdivision replaces a
vertex bound $b$ by $2^b$. Iterating this bound the supplied number of
times is primitive recursion. Carriers and their apices are tested by
exhaustive face scans. The row scans of elimination take $mn$ pivot
tests for insertion of $n$ columns, each using an exclusive-or on two
finite vectors. Integer vectors for the polygon calculation have arbitrary
signed entries as input. Every substitution is addition, multiplication,
or prefix summation on a bounded list, and extraction of the order-two
coordinate is division by two with remainder. The proof that equality of
coordinates entails the existence of a filling actually constructs that
filling by reversing the substitutions. Rational operations for the
polyhedron test use signed numerators, positive denominators, cross
multiplication and a test before division. Thus all these algorithms
are total after assigning a default on invalid input. Loops in the
proofs can use bounded Boolean predicates about the processed prefix,
including a finite table of the earlier partial answers to each face
recursion.

A Lean implementation can expose the following definitions in sequence.
Use \texttt{Nat} for masks, a dimension and a list of masks for columns,
and a pair of naturals or a signed constructor for integers. Give each
routine a numerical recursion parameter for its scan length.
\begin{enumerate}
\item Define \texttt{facesOfSize} and \texttt{deleteVertex}. Next build
\texttt{boundaryColumn}, \texttt{applyMatrix}, and
\texttt{composeMatrix}. Prove \texttt{boundarySquared} by pairing the
two orders of deleting distinct vertices. Implement
\texttt{reduceWithWitness} as downward pivot recursion, carrying the
preimage and the insertion index. Its two induction equalities say
that the stored column is the matrix applied to the stored preimage and
that processed pivot rows of the residual vanish. A separate prefix
invariant records triangularity of all processed domain vectors and
distinctness of nonzero pivots. These give \texttt{solve},
\texttt{kernelBasis}, and \texttt{homologyCoords} as above.
\item Define \texttt{coneColumn}, \texttt{simplicialImage}, and
\texttt{carrierHomotopy} on basis masks. Tabulate the last routine
in increasing face size. Its invariant is the checked matrix equality
$dP+Pd=F+G$ on all completed columns. Define \texttt{subdivision}
by the chain predicate on faces, and \texttt{subdivisionMap} by coning
the already available boundary. Apply the carrier lemma with the specific
cones in the subdivision and nerve proofs. The Mayer--Vietoris routine
uses Boolean assignment to split each face and \texttt{solve} on the
boundary columns for lifting.
\item Encode a signed letter as an edge identifier and a Boolean,
a walk as a start vertex and an edge list. Structural recursion gives
\texttt{reduceWord}. Prove its invariance under an inverse pair first on
one reduced stack, then under a common suffix. Encode \texttt{spokeBoundary}
as the signed sum of three indexed unit vectors. Check prefix substitution
and its inverse on a general signed input cycle, then specialize the
exponent vector by recursion along the blocks of the polygon word.
\item For the convex hull calculation implement \texttt{height},
\texttt{directionAt}, determinant signs and the two rational extrema
among pair interpolations. Build each side list through the extreme
ordinates at input abscissae, deleting collinear intermediate points.
The proof obligations are the three-weight support reduction, affine
interpolation inequalities, then cyclic block counting. All are finite
rational identities with sign tests. Check the global sum by counting
incidences at their unique highest vertex.
\item For barcodes store a list of matrices and dimensions and a list of
basis matrices embedding the current complements in the original
spaces. Compose rows backwards for $\lambda_t$, obtain its kernel by
\texttt{kernelBasis}, and express the restricted arrows in these
coordinates by \texttt{solve}. Recursion on the initial total dimension
as fuel terminates even if some stages are zero. Verify the direct-sum
identity stagewise and compute interval counts by four matrix ranks.
\end{enumerate}
